\documentclass{article}
\usepackage{graphicx}
\usepackage[en-GB]{datetime2}
\DTMlangsetup[en-GB]{ord=omit}

\title{Week 3 Example 4}
\author{Fernaldy}
\date{\today}

\begin{document}

\maketitle

\section*{ANOVA}
Scenario: A student trying to get some practice in carrying out ANOVA in SPSS decided to
compare weights (kg) between three independent groups (A, B \& C) with a null hypothesis
(H0) that there is no difference in mean weight between groups in the population.
The student first plotted histograms of weight for the 3 groups

\begin{figure}[htbp]
    \centering
    \includegraphics[width=\textwidth]{images/distribution.png}
    \caption{Histograms of weight for the 3 groups}
    \label{fig:histograms}
\end{figure}

\subsection*{Q1 - Comment on the distribution of weight in the 3 groups. Which assumption is being
checked here?}
The assumption being checked is for the normality in the distribution.
\newline
The distribution of weights in the three groups are roughly normal.

\subsection*{Q2a - What was the sample size? Comment on the size of the sample with respect to carrying
out ANOVA.}

\begin{figure}[htbp]
    \centering
    \includegraphics[width=\textwidth]{images/descriptives.png}
    \caption{Descriptive stats of weight for the 3 groups}
    \label{fig:summary table}
\end{figure}

The total sample size (all three groups) is 236, with
\begin{itemize}
    \item 81 samples from group A
    \item 88 samples from group B
    \item 67 samples from group C
\end{itemize}
Furthermore, the sample size is also roughly balanced and each group has more than enough sample.

\subsection*{Q2b - Which descriptives should be presented? Write them out}
The mean, sample size, and standard deviation should be presented for each group and the overall.

\begin{tabular}{|l|l|l|l|}
\hline
      & N   & Mean  & SD    \\ \hline
A     & 81  & 70.13 & 11.04 \\ \hline
B     & 88  & 69.28 & 13.03 \\ \hline
C     & 67  & 68.01 & 12.87 \\ \hline
Total & 236 & 69.21 & 12.31 \\ \hline
\end{tabular}
\newline

\subsection*{Q3a - What is the purpose of the test whose result is shown in the table below?}

\begin{figure}[htbp]
    \centering
    \includegraphics[width=\textwidth]{images/levene.png}
    \caption{Levene's Test}
    \label{fig:levene test}
\end{figure}

The Levene's test for homogeneity of variance is to check whether there is a significant
difference in variance between the sample groups and this test is used as part of ANOVA.


\subsection*{Q3b - Interpret the results of the table}
The value of \textit{Sig.} is 0.462. So, we can accept the null hypothesis and conclude
that there is \textbf{no significant difference} in the variances between the groups A, B, and C.
\\~\\
\begin{figure}[htbp]
    \centering
    \includegraphics[width=\textwidth]{images/anova.png}
    \caption{ANOVA results}
    \label{fig:anova}
\end{figure}

\textit{The student decided that the above table showed evidence that there is a difference in weight
between groups.}
\subsection*{Q4a - Is this statement correct?}
Incorrect, since the value of \textit{Sig.} is 0.579, we fail to reject the null hypothesis.
There is no
significant difference in the weight between the three groups.

\subsection*{Q4b - Interpret in full the results of the table.}
Since the null hypothesis is not rejected, we can conclude that there is no significant difference
in weight between the three groups.
\\~\\

\begin{figure}[htbp]
    \centering
    \includegraphics[width=\textwidth]{images/post-hoc.png}
    \caption{Post-hoc results}
    \label{fig:post-hoc}
\end{figure}

\subsection*{Q5 - What is your comment on this next step of the analysis? Was this an appropriate step
bearing in mind the ANOVA results?}
Since the result of the ANOVA test showed that there is no significant difference in weight
between the three groups $p > 0.05$ the post-hoc test is redundant and inappropriate.


\end{document}
